% Annexe du bloc « fiscalité du capital », révision v3 (note8p/CONSIGNES_V3.md). Chiffres : note8p/sorties/pfu3.json et
% pfu3_calculs.txt (programme note8p/calculs/pfu3_neutralite.py) ; coût d'un point de PFU, CEHR et CDHR :
% note8p/sorties/pfu.json (programme note8p/calculs/pfu_transmission.py). Figure : note8p/figures/g9_exemple_pv.

\subsection{La règle : reprise du prix, correction à la marge}

Le PFU reste à 31,4~\% : 12,8~\% d'IR (CGI, art.~200~A) et 18,6~\% de prélèvements sociaux (LFSS pour 2026,
art.~12). À la donation et au décès, la plus-value des titres de l'art.~150-0~A n'est plus effacée : le bénéficiaire
reprend le prix d'acquisition historique et n'est imposé que s'il vend [C]. Reste à éviter que l'impôt latent paie
aussi des droits de mutation. Notons $V$ la valeur transmise, $G$ la plus-value, $t$ le PFU et $D(\cdot)$ les droits
du bénéficiaire sur sa part. Le parent qui vend paie $tG + D(V-tG)$ ; l'héritier qui reçoit puis vend paie $D(V) +
t(G-a)$, où $a$ est la majoration du prix repris. Les deux charges sont égales si et seulement si
\[
t\,a \;=\; D(V) - D(V - tG),
\]
c'est-à-dire si l'héritier récupère, à la vente, les droits qu'il a payés sur l'impôt latent. Les droits ne doivent
pas frapper la part de l'héritage qui reviendra au fisc. D'où la règle retenue [C], dite « correction à la marge » :
\emph{le prix repris est majoré de $a = [D(V) - D(V - tG)]/t$}. Le notaire calcule deux fois les droits du bénéficiaire,
une fois sur sa part, une fois sur sa part diminuée de l'impôt latent. Il inscrit la différence dans la déclaration,
qui doit de toute façon porter le prix historique. La majoration s'attache aux titres : elle suit les ventes
partielles au prorata et se cumule si les titres sont transmis à nouveau. Elle équivaut à un crédit d'impôt à la vente,
mais sous la forme d'un prix, plus simple à suivre.

\begin{table}[!htbp]\centering\scriptsize
\caption{Règles comparées}\label{tab:cap_regles}
\begin{tabular}{@{}p{4.4cm}p{3.1cm}p{3.2cm}p{3.2cm}@{}}
\toprule
Règle & Neutralité vendre / transmettre & Si l'héritier ne vend jamais & Coût et simplicité \\
\midrule
(i) Taux moyen des droits $\times G$ ajouté au prix (v2 ; 26 U.S.C. §~1015(d)(6)) & Double imposition résiduelle
de $tG(m-\bar\tau)$ : 6~\% de la charge à 1~M€ & Droits inchangés & Simple ; recette de référence \\
(ii) Même correction au taux marginal $m$ & Exacte tant que $V$ et $V-tG$ sont dans la même tranche & Droits inchangés
& Simple ; 0,2 à 0,3~Md€/an de moins que (i) \\
(iii) Impôt latent déduit de l'assiette des droits & Exacte & Droits \emph{réduits} pour un impôt jamais payé &
Coûte 0,8 à 1,1~Md€/an dès 2027 \\
\textbf{(iv) Crédit à la vente des droits afférents à l'impôt latent (retenue, sous forme de prix)} & Exacte & Droits
inchangés & Un second calcul de droits ; comme (ii) \\
(v) Prix majoré de 30~\% de $G$ & Juste si $m = 30~\%$ ; sinon jusqu'à $-30~\%$ (conjoint) ou $+10~\%$ (neveu) &
Droits inchangés & Simple mais arbitraire \\
(vi) Réalisation au décès, impôt déduit des droits (Canada) & Exacte & Impôt dû quand même & Contraire à
l'imposition à la vente seule (non retenue) \\
\bottomrule
\end{tabular}
\end{table}

\begin{table}[!htbp]\centering\scriptsize\setlength{\tabcolsep}{4pt}
\caption{Charge totale (plus-value + droits) et écart avec « vendre puis transmettre »}\label{tab:cap_cas}
\begin{tabular}{@{}lrrrrrr@{}}
\toprule
 & Vendre, puis & Transmettre & \multicolumn{4}{c}{Réforme : écart (transmettre puis vendre)} \\
\cmidrule(l){4-7}
Cas (droit de 2026) & transmettre & aujourd'hui & (i) moyen & (iv) retenue & (iii) & (v) 30~\% \\
\midrule
100~k€ devenus 1~M€, un enfant ($\bar\tau$ 21~\%, $m$ 30~\%) & 410\,782 & 212\,962 & +24\,597 & 0 & 0 & 0 \\
250~k€ devenus 400~k€, un enfant (15~\%, 20~\%) & 95\,874 & 58\,194 & +2\,568 & 0 & 0 & $-$4\,710 \\
2~M€ devenus 40~M€, deux enfants (44~\%, 45~\%) & 23\,997\,389 & 17\,434\,789 & +168\,603 & 0 & 0 & +1\,789\,800 \\
100~k€ devenus 1~M€, un frère (44~\%, 45~\%) & 595\,818 & 440\,388 & +2\,716 & 0 & 0 & +42\,390 \\
100~k€ devenus 1~M€, un neveu (55~\%, 55~\%) & 672\,788 & 545\,618 & +1\,238 & 0 & 0 & +70\,650 \\
100~k€ devenus 1~M€, conjoint (exonéré) & 282\,600 & 0 & 0 & 0 & 0 & $-$84\,780 \\
\bottomrule
\end{tabular}
\par\vspace{3pt}\parbox{0.95\linewidth}{\raggedright\scriptsize Lecture : aujourd'hui, transmettre avant de vendre
divise la charge par deux dans l'exemple ; avec la règle (iv), les deux chemins coûtent autant. L'écart de la règle (i)
vient de la progressivité et de l'abattement : il disparaît presque aux taux plats (neveu) ou au taux maximal
(très grande fortune). Si l'héritier ne vend jamais, il paie avec (i), (ii), (iv) et (v) les droits d'aujourd'hui ;
avec (iii), moins qu'aujourd'hui (128\,182~€ au lieu de 212\,962~€ dans l'exemple).}
\end{table}

\subsection{Recettes 2027-2040 de la règle retenue}

Méthode, jeux d'hypothèses et calendrier sont ceux de la version 2 (\texttt{pfu\_transmission.py}) : impôt latent
transmis par an $L = 31{,}4~\%\,E\,b\,\lambda + 30~\%\,S/T$ (3,65 / 7,49 / 14,68~Md€ dans les jeux bas, central,
haut), part $a$ vendue dans les douze mois puis $h$ par an, impôt encaissé l'année qui suit la vente. Seul change le
facteur de correction : $(1-m)$ au lieu de $(1-\bar\tau)$, avec $\bar\tau$ = 15 / 10 / 5~\% [H] et
$m = 1{,}4\,\bar\tau$ [H]. Le rapport 1,4 est celui des cas-types en ligne directe (1,0 au taux maximal, 1,4 à 1~M€, 1,6
à 1,8 sous pacte Dutreil). Unités : Md€ au niveau du PIB de 2025.

\begin{table}[!htbp]\centering\scriptsize\setlength{\tabcolsep}{3pt}
\caption{Recette encaissée par année (Md€)}\label{tab:cap_annuel}
\begin{tabular}{@{}l*{14}{r}r@{}}
\toprule
 & 27 & 28 & 29 & 30 & 31 & 32 & 33 & 34 & 35 & 36 & 37 & 38 & 39 & 40 & Cumul \\
\midrule
\multicolumn{16}{@{}l}{\emph{Règle retenue (iv)}} \\
Bas & 0 & 0,19 & 0,44 & 0,49 & 0,54 & 0,58 & 0,63 & 0,67 & 0,72 & 0,76 & 0,80 & 0,85 & 0,89 & 0,93 & 8,5 \\
Central & 0 & 0,82 & 1,83 & 2,01 & 2,19 & 2,36 & 2,52 & 2,68 & 2,83 & 2,97 & 3,11 & 3,24 & 3,37 & 3,50 & 33,4 \\
Haut & 0 & 2,51 & 5,53 & 6,02 & 6,48 & 6,91 & 7,31 & 7,69 & 8,05 & 8,38 & 8,70 & 9,00 & 9,28 & 9,54 & 95,4 \\
\textbf{Prudent} (bas, deux ans plus tard) & 0 & 0 & 0 & 0,19 & 0,44 & 0,49 & 0,54 & 0,58 & 0,63 & 0,67 & 0,72 & 0,76 & 0,80 & 0,85 & 6,7 \\
\multicolumn{16}{@{}l}{\emph{Pour comparaison, jeu central}} \\
(i) taux moyen (v2) & 0 & 0,86 & 1,91 & 2,10 & 2,29 & 2,47 & 2,64 & 2,80 & 2,96 & 3,11 & 3,26 & 3,40 & 3,53 & 3,66 & 35,0 \\
(iii) déduction des droits & $-$1,05 & $-$0,10 & 1,08 & 1,29 & 1,50 & 1,69 & 1,88 & 2,07 & 2,24 & 2,41 & 2,57 & 2,72 & 2,87 & 3,02 & 24,2 \\
\bottomrule
\end{tabular}
\par\vspace{3pt}\parbox{0.95\linewidth}{\raggedright\scriptsize La neutralité coûte 1,6~Md€ cumulés sur 2027-2040 au
centre (0,6 à 2,1), soit 0,3~Md€ par an en régime de croisière. La règle (iii) coûte dès 2027, avant toute recette, et
donne aux héritiers qui ne vendent pas une baisse des droits ; sur le jeu bas, elle ne rapporte rien en cumul. La
règle (vi) rapporterait 6,4~Md€ par an dès 2028 au centre, mais taxe à la transmission. Le scénario prudent ne sert
qu'au bouclage : il reprend le jeu bas, avec deux ans de décalage, pour n'engager que des recettes constatées.}
\end{table}

\subsection{Pacte Dutreil, localisation des grands patrimoines, exit tax}

\emph{Pacte Dutreil} (art.~787~B : 75~\% de la valeur exonérée ; engagement individuel porté de 4 à 6~ans par la
LF 2026, sources secondaires [S]). La plus-value latente suit les titres et la correction porte sur les droits
réellement payés, qui sont faibles : $m$ vaut 5 à 10~\%. Entreprise de 10~M€ (prix 0,5~M€) transmise à deux enfants,
qui la cèdent après l'engagement : charge totale de 0,63~M€ aujourd'hui (6~\%), 3,34~M€ avec la réforme (33~\%), contre
5,58~M€ (56~\%) si le fondateur avait vendu avant de transmettre. Dans l'exemple à 1~M€ : 3~\% aujourd'hui, 30~\% avec
la réforme, 41~\% en vendant avant. L'écart restant n'est pas une interaction des deux impôts : c'est l'exonération
Dutreil elle-même, conservée par l'héritier qui vend après l'engagement. Si les héritiers gardent l'entreprise, rien ne
change. La hausse de charge pour les familles qui cèdent est forte, et elle est voulue [C] : le cumul « 75~\%
d'exonération + effacement » est supprimé, pas le pacte.

\emph{Localisation.} Une fois l'effacement fermé, la sortie qui reste est le départ. L'exit tax (art.~167~bis) impose
la plus-value latente au départ (participation d'au moins 50~\% ou titres de plus de 800\,000~€), avec sursis, puis
dégrève l'impôt si les titres sont encore détenus deux ans après le départ, cinq ans au-delà de 2,57~M€ [S : sources
secondaires, texte non relu]. La plus-value reprise entre dans sa base. Un héritier déjà non résident qui vend des
titres de portefeuille échappe en général à l'impôt français (conventions non vérifiées). La littérature dit trois
choses [E]. (1) Les très riches bougent : un point d'impôt sur la fortune en plus réduit de 2~\% environ le nombre de
contribuables fortunés en Scandinavie (Jakobsen, Kleven, Kolsrud, Landais et Muñoz, NBER WP 32153, \emph{AER} 2026) ;
après 2011, Madrid, seule région espagnole sans impôt sur la fortune, gagne 7,5~\% de contribuables fortunés en six ans
(Agrawal, Foremny et Martínez-Toledano, \emph{AEJ: Applied} 2025) ; 35~\% des milliardaires du Forbes 400 quittent
les États américains qui ont gardé un impôt successoral après 2001 (Moretti et Wilson, \emph{AEJ: Policy} 2023).
(2) Les recettes résistent : à Madrid comme aux États-Unis, l'impôt rapporte plus qu'il ne fait perdre, et en
Scandinavie 0,22~\$ seulement sur chaque dollar levé est perdu par départs. (3) En France, la réaction mesurée est
faible au niveau du dernier centile (0,02 à 0,23~\% des foyers de plus par point d'imposition effective, CAE,
Focus n°~118, 2025) ; les retours de redevables de l'IFI dépassent les départs depuis 2018 (France Stratégie, 2023).
Rien ne mesure le comportement du dernier dix-millième, qui porte l'essentiel de l'impôt latent. Ces élasticités
portent sur des impôts annuels sur le stock et ne se transposent pas telles quelles à un impôt payé une fois, à la
vente. Nous ne les appliquons pas [choix de méthode] : la prudence passe par le scénario de recettes prudent et par
deux mesures de garde [C] : délai de dégrèvement de l'exit tax porté à quinze ans pour les participations de plus de
2,57~M€ (durée en vigueur de 2014 à 2018, sources secondaires), et imposition, dans les limites des conventions, de
la plus-value transmise à un bénéficiaire non résident (à expertiser).

\subsection{Taux du PFU, CEHR et CDHR (inchangés depuis la version 2)}

Un point de PFU vaut 0,67~Md€ à comportements inchangés sur l'assiette des revenus de 2025 (67,0~Md€), 0,44~Md€ les
premières années avec l'élasticité de 0,77 de Lefebvre, Lehmann et Sicsic (2025) [S]. Une baisse à 25~\% coûterait
4,3~Md€ en mécanique, 3,1~Md€ après réaction. Les évaluations du PFU de 2018 ne trouvent d'effet que sur les dividendes
distribués (+15~\%, note IPP n°~46, 2019), aucun sur l'investissement (France Stratégie, 2023 ; Yagan, 2015, \emph{AER}) ;
Boissel et Matray (2022) ont été rétractés (\emph{AER} 113(7), 2023). CEHR et CDHR sont maintenues sur le capital
(coût nul) ; les supprimer coûterait 0,7 à 1,1 et 0,6~Md€. La CRDS s'éteint fin 2033 : le PFU passera à 30,9~\% sans
mesure nouvelle. L'indexation des plus-values sur l'inflation (0,6 à 0,9~Md€ par an) reste une orientation de long
terme.
